% ============================================================
%  SASTO ISEF Poster v5  —  48" × 36"  (121.92 × 91.44 cm)
%  Three panels: Left 12" | Center 24" | Right 12"
%
%  Compile:  pdflatex poster_v5.tex
%            (run twice for correct layout)
%
%  All images use paths relative to this file:
%    poster_final/renders_hq/   <- new HQ 3D renders
%    poster_final/renders/      <- per-geometry renders
%    poster_figures_v5/         <- matplotlib data figures
%    poster_images_extracted/   <- extracted diagram images
% ============================================================

\documentclass[10pt]{article}

\usepackage[
  paperwidth=121.92cm, paperheight=91.44cm,
  top=0cm, bottom=0cm, left=0cm, right=0cm,
  noheadfoot
]{geometry}

\usepackage[T1]{fontenc}
\usepackage[utf8]{inputenc}
\usepackage{lmodern}
\usepackage{microtype}
\usepackage{amsmath,amssymb}
\usepackage{graphicx}
\usepackage{xcolor}
\usepackage[most]{tcolorbox}
\usepackage{enumitem}
\usepackage{booktabs}
\usepackage{colortbl}
\usepackage{tabularx}
\usepackage{array}
\usepackage{multirow}
\usepackage{calc}
\usepackage{tikz}
\usepackage{ragged2e}
\usepackage{setspace}

\pagestyle{empty}
\setlength{\parindent}{0pt}
\setlength{\parskip}{0pt}

% ── COLOURS ─────────────────────────────────────────────────
\definecolor{bgnavy}   {HTML}{062B7A}
\definecolor{titleband}{HTML}{032061}
\definecolor{secbar}   {HTML}{0A3D9A}
\definecolor{cardfill} {HTML}{F7F9FC}
\definecolor{cardbdr}  {HTML}{B7C5E3}
\definecolor{txtdark}  {HTML}{0B1736}
\definecolor{posterred}{HTML}{D7263D}
\definecolor{teal}     {HTML}{008C9E}
\definecolor{gold}     {HTML}{CFA535}
\definecolor{eqbg}     {HTML}{E8EEF2}
\definecolor{altrow}   {HTML}{EEF2FA}
\definecolor{darkcard} {HTML}{1A3A6E}

% ── CONTENT CARD ─────────────────────────────────────────────
\newenvironment{ccard}{%
  \begin{tcolorbox}[
    colback=cardfill, colframe=cardbdr,
    boxrule=0.5pt, arc=4pt, outer arc=4pt,
    left=3.5mm, right=3.5mm, top=2.5mm, bottom=2.5mm,
    width=\linewidth
  ]%
}{%
  \end{tcolorbox}%
}

% ── SECTION HEADER ───────────────────────────────────────────
\newcommand{\sechdr}[1]{%
  \begin{tcolorbox}[
    colback=secbar, colframe=secbar,
    boxrule=0pt, arc=3pt,
    left=4mm, right=4mm, top=1mm, bottom=1mm,
    width=\linewidth
  ]
    \centering\color{white}\fontsize{24pt}{26pt}\selectfont\bfseries
    \MakeUppercase{#1}
  \end{tcolorbox}%
  \vspace{1.5mm}%
}

% ── SUBSECTION LABEL ─────────────────────────────────────────
\newcommand{\sub}[1]{%
  {\color{secbar}\bfseries\large #1}\par\vspace{1mm}}

% ── EQUATION PILL ────────────────────────────────────────────
\newcommand{\eqpill}[1]{%
  \begin{tcolorbox}[
    colback=eqbg, colframe=cardbdr,
    boxrule=0.5pt, arc=5pt,
    left=3mm, right=3mm, top=1mm, bottom=1mm,
    width=\linewidth
  ]
    \centering\small #1
  \end{tcolorbox}%
}

% ── CRITERIA BOX ─────────────────────────────────────────────
\newcommand{\cbox}[2]{%
  \begin{tcolorbox}[
    colback=darkcard, colframe=cardbdr,
    boxrule=0.5pt, arc=5pt,
    left=2mm, right=2mm, top=2mm, bottom=2mm,
    width=\linewidth
  ]
    \centering{\color{gold}\bfseries\large #1}\\[3pt]
    \color{white}\small #2
  \end{tcolorbox}%
}

% ── FIGURE CAPTION ───────────────────────────────────────────
\newcommand{\fc}[1]{{\centering\color{txtdark}\tiny\itshape #1\par}\vspace{1mm}}

% ── IMAGE COMMAND ────────────────────────────────────────────
\newcommand{\fig}[2]{\includegraphics[width=#2,keepaspectratio]{#1}}

% ── PATH SHORTCUTS ───────────────────────────────────────────
\newcommand{\HQ}[1]{poster_final/renders_hq/#1}
\newcommand{\RD}[1]{poster_final/renders/#1}
\newcommand{\PF}[1]{poster_figures_v5/#1}
\newcommand{\EX}[1]{poster_images_extracted/#1}

% ── DIMENSIONS ───────────────────────────────────────────────
% Poster = 121.92cm wide; panels: 30.48 | 60.96 | 30.48 cm
% Inner margin: 0.38cm each side
\newlength{\LC}  \setlength{\LC} {29.72cm}   % left col content width
\newlength{\CC}  \setlength{\CC} {60.20cm}   % center col content width
\newlength{\RC}  \setlength{\RC} {29.72cm}   % right col content width
\newlength{\HC}  \setlength{\HC} {29.60cm}   % half of center col
\newlength{\TC}  \setlength{\TC} {19.33cm}   % third of center col
\newlength{\SG}  \setlength{\SG} {1.5mm}     % section gap
\newlength{\PM}  \setlength{\PM} {0.38cm}    % panel margin

% ============================================================
\begin{document}
\noindent
% ============================================================
%  TITLE BAND  (3" = 7.62cm, full width)
% ============================================================
\begin{tikzpicture}[remember picture,overlay]
  \fill[titleband] (current page.north west) rectangle ++(121.92cm,-7.62cm);
  \fill[gold] ([yshift=-7.55cm]current page.north west) rectangle ++(121.92cm,-0.07cm);
\end{tikzpicture}%
\begin{minipage}[t][7.62cm][c]{121.92cm}
  \hspace{0.5cm}%
  % Left 3D renders
  \begin{minipage}[c]{6.4cm}
    \fig{\RD{thumb_original_house.png}}{6.4cm}\\[-1pt]
    {\centering\tiny\color{white}Original structure\par}
  \end{minipage}%
  \hspace{0.25cm}%
  \begin{minipage}[c]{6.4cm}
    \fig{\RD{thumb_optimized_house.png}}{6.4cm}\\[-1pt]
    {\centering\tiny\color{white}After SASTO-PA\par}
  \end{minipage}%
  \hfill
  % Center title block
  \begin{minipage}[c]{76cm}
    \centering\color{white}
    {\fontsize{46pt}{48pt}\selectfont\bfseries
      SURROGATE-ACCELERATED STRUCTURAL TOPOLOGY OPTIMIZATION}\\[4pt]
    {\fontsize{23pt}{25pt}\selectfont\bfseries\itshape
      A Deep Ensemble Surrogate for Fast, Printable, and Safety-Verified Building Optimization}\\[5pt]
    {\fontsize{20pt}{22pt}\selectfont\bfseries Eric Hou
      \hspace{1cm}{\normalfont\fontsize{15pt}{17pt}\selectfont
        Engineering Category \hspace{5mm} SASTO Project \hspace{5mm} March 2026}}
  \end{minipage}%
  \hfill
  % Right credit block
  \begin{minipage}[c]{22cm}
    \RaggedRight\small\color{white}
    {\bfseries\large Credit Line of Origin}\\[3pt]
    {\bfseries Credit Line of Origin}\\[4pt]
    {\itshape References \& Data: see Key References.}\\[4pt]
    * Images with asterisk are public domain\\
    (Wikimedia Commons CC-BY-SA~3.0).\\[3pt]
    All other graphics and images created by\\
    Eric Hou, 2026, unless otherwise noted.
  \end{minipage}%
  \hspace{0.38cm}
\end{minipage}

\vspace{2mm}

% ============================================================
%  THREE PANELS
% ============================================================
\noindent\hspace{\PM}%
%% ════════════════════════════════════════════════════════════
%%  LEFT PANEL  (29.72cm)
%% ════════════════════════════════════════════════════════════
\begin{minipage}[t]{\LC}

%% L1: VISUAL ABSTRACT ─────────────────────────────────────
\sechdr{Visual Abstract}
\begin{ccard}
  \centering
  \fig{\HQ{fig_wireframe_pipeline.png}}{\linewidth}\\[2pt]
  \fc{Fig.~1.\ Full SASTO pipeline:
    3DWire wireframe $\to$ volumetric extrusion $\to$ FEA simulation $\to$
    voxel encoding $\to$ deep ensemble training $\to$ sensitivity-guided erosion $\to$
    watertight optimized STL.}

  \RaggedRight\small\vspace{1mm}
  \textbf{a)~User Workflow:}
  Given a building wireframe, SASTO extrudes walls, roof, and floor into a
  watertight volumetric solid; runs a single FEA simulation; then applies
  surrogate-guided erosion in 50\,s to produce an optimized single-component STL
  with 23--92$\times$ speedup over SIMP.\\[3pt]
  \textbf{b)~Model Creation:}
  14,293 wireframes from the 3DWire dataset were processed into volumetric meshes.
  After filtering 3,115 diverged/invalid solves, \textbf{11,178 valid 128-cube voxel
  grids} remain, split family-aware into \textbf{8,943 train / 1,121 val / 1,114 test}.
  Each sample is labeled with peak Von Mises stress, max displacement, and global compliance.
\end{ccard}
\vspace{\SG}

%% L2: INTRODUCTION ────────────────────────────────────────
\sechdr{Introduction}
\begin{ccard}
  % ── Opening prose + right-floated image (mirrors reference poster layout) ──
  \begin{minipage}[t]{18.5cm}
    \RaggedRight\small
    \sub{CONCRETE \& GLOBAL EMISSIONS}
    Cement production accounts for approximately \textbf{8\% of global CO$_2$ emissions}
    [IEA~2021], more than all commercial aviation combined.
    A single residential building consumes 30--80 tonnes of concrete, most placed in walls
    sized for the \textit{worst-case load scenario} rather than the service condition.
    This uniform-thickness paradigm is an artifact of conventional formwork, not structural
    optimality.\\[4pt]
    \textbf{Large-scale 3D concrete printing} companies (ICON, COBOD, Apis~Cor) can
    now extrude \textit{arbitrary wall-thickness profiles} at no additional tooling cost.
    This makes topology-optimized heterogeneous walls \textbf{practically realizable at
    building scale for the first time.}
  \end{minipage}\hfill
  \begin{minipage}[t]{9.8cm}\centering
    \fig{\HQ{wireframe_panel.png}}{\linewidth}\\[-2pt]
    \fc{3DWire building wireframes (Lin et al.\ 2024)}
    \vspace{1.5mm}
    \fig{\RD{REF_original_colored_iso.png}}{\linewidth}\\[-2pt]
    \fc{Volumetric mesh (reference case)}
  \end{minipage}

  % ── 3-box causal chain ──────────────────────────────────────────────────────
  \vspace{3mm}
  \noindent
  \begin{minipage}[c]{8.2cm}
    \begin{tcolorbox}[colback=secbar!12,colframe=secbar,arc=3pt,
      left=2mm,right=2mm,top=1.5mm,bottom=1.5mm,width=\linewidth]
      \centering\small\bfseries\color{secbar}
      Worst-case wall sizing\\[-1pt]
      {\normalfont\footnotesize wastes 30--80\,t concrete per building}
    \end{tcolorbox}
  \end{minipage}%
  \hfill{\color{gold}\Large$\Rightarrow$}\hfill
  \begin{minipage}[c]{8.8cm}
    \begin{tcolorbox}[colback=secbar!12,colframe=secbar,arc=3pt,
      left=2mm,right=2mm,top=1.5mm,bottom=1.5mm,width=\linewidth]
      \centering\small\bfseries\color{secbar}
      FEA ``hidden'' cost\\[-1pt]
      {\normalfont\footnotesize 100$+$ solves $\times$ 19--77\,min at $128^3$}
    \end{tcolorbox}
  \end{minipage}%
  \hfill{\color{gold}\Large$\Rightarrow$}\hfill
  \begin{minipage}[c]{9.0cm}
    \begin{tcolorbox}[colback=secbar!12,colframe=secbar,arc=3pt,
      left=2mm,right=2mm,top=1.5mm,bottom=1.5mm,width=\linewidth]
      \centering\small\bfseries\color{secbar}
      Unprintable geometry\\[-1pt]
      {\normalfont\footnotesize diagonal voxels $\to$ floating fragments}
    \end{tcolorbox}
  \end{minipage}

  % ── Subsection pivot + equation with right-side annotation ──────────────────
  \vspace{3mm}
  \sub{SURROGATE TOPOLOGY OPTIMIZATION}
  \small\RaggedRight
  Classical \textbf{SIMP} topology optimization iterates 100--1000 full FEA solves.
  \textbf{SASTO} replaces this loop with a trained surrogate, yet existing approaches
  are not used in building-scale design because:

  \vspace{2mm}
  \begin{minipage}[c]{0.44\linewidth}
    \centering
    \[{\color{secbar}\large\mathbf{K\,u = f}}\]
  \end{minipage}%
  \begin{minipage}[c]{0.54\linewidth}
    \small\RaggedRight
    \textbf{Equation:} Linear elastostatics.\\
    $K$: global stiffness matrix \hspace{1mm} $u$: displacement field\\
    $f$: applied load vector.\\[2pt]
    Solved \textbf{once per design} during data generation;
    \textbf{never} during optimization in SASTO vs.\ 100$+$ times per design in SIMP.
  \end{minipage}

  \vspace{2mm}
  \begin{enumerate}[leftmargin=1.6em, itemsep=2pt, topsep=0pt, label=\textbf{(\arabic*)}]
    \item FEA requires 100$+$ expensive iterative solves, making it infeasible at building
      scale ($128^3$ mesh, 19--77\,min per solve).
    \item No guarantee of topological connectivity for \textbf{3D-printable
      single-component toolpaths}: diagonal voxel adjacency produces
      unfabricable floating geometry.
    \item Uniform wall-thickness assumption \textbf{wastes material} in structurally
      unloaded interior partitions.
  \end{enumerate}

  \vspace{2mm}
  \small\RaggedRight
  \textbf{SASTO} addresses all three simultaneously: a trained surrogate replaces
  iterative FEA, 6-connectivity enforcement guarantees single-component printable
  geometry, and part-aware thickness floors maximize material savings where structural
  demand is lowest.
\end{ccard}
\vspace{\SG}

%% L3: RESEARCH OBJECTIVES ──────────────────────────────────
\sechdr{Research Objectives}
\begin{ccard}
  \small\RaggedRight
  \textbf{ENGINEERING PROBLEM:} SIMP topology optimization requires 100s--1000s of costly
  FEA solves; existing surrogate models lack formal topology guarantees; and uniform-thickness
  constraints forego material savings in structurally non-critical interior walls.\\[4pt]
  \textbf{RESEARCH QUESTION:} Can a deep learning surrogate replace iterative FEA to
  enable fast, 3D-printable, and safety-verified topology optimization of full building
  structures, validated against independent ground-truth FEA on 1,114 held-out designs?\\[4pt]
  \textbf{CLAIM:} \textbf{SASTO} replaces iterative FEA with a 5-member deep ensemble
  providing uncertainty-quantified structural predictions. It enforces a \textbf{6-connected
  topology} criterion for printability and achieves \textbf{23--92$\times$ speedup} with
  \textbf{zero safety violations} on 1,114 test designs, verified by independent hex8 FEA.

  \vspace{3mm}
  \begin{minipage}[t]{0.31\linewidth}
    \cbox{Speed}{50\,s median runtime.\\23--92$\times$ vs.\ SIMP.}
  \end{minipage}\hfill
  \begin{minipage}[t]{0.31\linewidth}
    \cbox{Printability}{6-connectivity enforced.\\Single-component STL.}
  \end{minipage}\hfill
  \begin{minipage}[t]{0.31\linewidth}
    \cbox{Safety}{$k$=1.0 bound.\\0/1,114 FEA violations.}
  \end{minipage}

  \vspace{3mm}
  \textbf{Primary Endpoint:} Minimize material volume subject to structural constraints
  (stress $\leq \sigma_\text{allow}$, compliance ratio $\leq 1.15$) with zero violations
  confirmed by independent hex8 FEA across the full 1,114-geometry test set.\\[3pt]
  \textbf{Secondary Endpoint:} Compliance Spearman $\rho \geq 0.90$; $\geq 90\%$ single-component
  STL meshes without post-processing.\\[4pt]
  \textbf{BRAINSTORMING A SOLUTION:}\\
  \begin{itemize}[leftmargin=1.2em, itemsep=1pt, topsep=0pt]
    \item \textbf{Data:} Wireframe-to-FEA pipeline using 3DWire + SfePy; 14,293 samples.
    \item \textbf{Surrogate:} 5 independently seeded 3D-CNNs; ensemble mean $+$ s.d.\ as
      uncertainty-quantified predictions of stress, displacement, and compliance.
    \item \textbf{Optimize:} Per-voxel sensitivity guides removal; 6-simple-point test guards
      connectivity; trust-region adaptive batch avoids constraint overshoot.
    \item \textbf{Validate:} Re-run independent ground-truth FEA on every optimized design;
      compute conformal coverage certificates.
  \end{itemize}
\end{ccard}
\vspace{\SG}

%% L4: PROBLEM FRAMING ─────────────────────────────────────
\sechdr{Problem Framing}
\begin{ccard}
  \small\RaggedRight
  \sub{Optimization Objective}
  The three-term objective balances volume removal, surface smoothness, and structural
  constraint satisfaction. Both volume $V$ and surface area $S$ are normalized by prior
  volume $V_0$; the constraint penalty $P$ aggregates scaled violations:
  \[\min_{\mathbf{p}\in\{0,1\}^N}\;
    J(\mathbf{p}) = w_V\!\left(\!\frac{V}{V_0}\!\right)
    + w_S\!\left(\!\frac{S}{V_0}\!\right) + P_\text{constraint}(\mathbf{p})\]
  Weights $w_V\!=\!1.0$, $w_S\!=\!0.01$ are selected by grid search on the validation
  split.  $P$ aggregates stress, compliance, and displacement violations, each normalized
  to a $[0,1]$ penalty scale.

  \vspace{2mm}
  \sub{Sensitivity Computation}
  Each iteration, back-propagation through the ensemble yields a per-voxel removal score:
  \[s_i = \frac{1}{M}\sum_{m=1}^{M}\frac{\partial}{\partial p_i}
    \bigl[f_m(C) + 0.3\,f_m(\sigma)\bigr]\]
  {\color{teal}$s_i > 0$: voxel penalizes objectives more than it stiffens; safe to remove.}\\
  {\color{posterred}$s_i < 0$: voxel is load-bearing; removal risks violation.}

  \vspace{2mm}
  \sub{Part-Aware Thickness (Cross-Section Views)}
  \begin{minipage}[c]{0.31\linewidth}\centering
    \fig{\HQ{xs_original.png}}{8.6cm}\\[-1pt]\fc{Original}
  \end{minipage}\hfill
  \begin{minipage}[c]{0.31\linewidth}\centering
    \fig{\HQ{xs_sasto_u.png}}{8.6cm}\\[-1pt]\fc{SASTO-U}
  \end{minipage}\hfill
  \begin{minipage}[c]{0.31\linewidth}\centering
    \fig{\HQ{xs_sasto_pa.png}}{8.6cm}\\[-1pt]\fc{SASTO-PA}
  \end{minipage}

  \vspace{2mm}
  Exterior walls, roof, and floor require $t_{\min}=2\Delta x\approx 156\,\text{mm}$ for
  structural and weather-tightness integrity. Interior partitions carry negligible vertical
  load and are permitted to thin to $t_{\min}=1\Delta x\approx 78\,\text{mm}$.
  This \textbf{part-aware} policy yields an additional \textbf{10.7\,pp} volume reduction
  vs.\ the uniform variant (SASTO-U: 34.3\% $\to$ SASTO-PA: 45.0\% on the reference case),
  as visible in the cross-section comparison above.
\end{ccard}

\end{minipage}%  ─── END LEFT PANEL ───
\hspace{3mm}%
%% ════════════════════════════════════════════════════════════
%%  CENTER PANEL  (60.20cm)
%% ════════════════════════════════════════════════════════════
\begin{minipage}[t]{\CC}

%% C1: ENGINEERING METHODOLOGY ──────────────────────────────
\sechdr{Engineering Methodology}
\begin{ccard}

% Overview + step pills row
\begin{minipage}[t]{30cm}
  \small\RaggedRight
  \textbf{Overview:}
  My methodology has four stages: (1)~dataset generation from real building wireframes via
  FEA, (2)~3D-CNN ensemble training with uncertainty quantification, (3)~SASTO
  sensitivity-guided voxel erosion with topology enforcement, and (4)~independent FEA
  ground-truth validation on all 1,114 held-out designs.
\end{minipage}\hfill
\begin{minipage}[t]{28cm}\centering
  \begin{tcolorbox}[colback=teal,colframe=teal,arc=3pt,left=1mm,right=1mm,
    top=0.5mm,bottom=0.5mm,width=5.8cm,nobeforeafter]
    \centering\color{white}\bfseries\footnotesize Wireframe\\\itshape Step 1
  \end{tcolorbox}\hspace{0.5mm}%
  {\color{gold}\small$\to$}\hspace{0.5mm}%
  \begin{tcolorbox}[colback=teal,colframe=teal,arc=3pt,left=1mm,right=1mm,
    top=0.5mm,bottom=0.5mm,width=5.8cm,nobeforeafter]
    \centering\color{white}\bfseries\footnotesize Volumetric\\\itshape Step 2
  \end{tcolorbox}\hspace{0.5mm}%
  {\color{gold}\small$\to$}\hspace{0.5mm}%
  \begin{tcolorbox}[colback=teal,colframe=teal,arc=3pt,left=1mm,right=1mm,
    top=0.5mm,bottom=0.5mm,width=5.8cm,nobeforeafter]
    \centering\color{white}\bfseries\footnotesize FEA + Train\\\itshape Step 3
  \end{tcolorbox}\hspace{0.5mm}%
  {\color{gold}\small$\to$}\hspace{0.5mm}%
  \begin{tcolorbox}[colback=teal,colframe=teal,arc=3pt,left=1mm,right=1mm,
    top=0.5mm,bottom=0.5mm,width=5.8cm,nobeforeafter]
    \centering\color{white}\bfseries\footnotesize Optimize\\\itshape Step 4
  \end{tcolorbox}
\end{minipage}

\vspace{3mm}\noindent\rule{\linewidth}{0.3pt}\vspace{2mm}

% Data Generation | Deep Ensemble
\begin{minipage}[t]{\HC}
  \sub{(1) Data Generation Pipeline}
  \fig{\HQ{fig_wireframe_pipeline.png}}{\linewidth}\\[-1pt]
  \fc{Fig.~2.\ Pipeline: 3DWire wireframe $\to$ boolean-fused mesh $\to$ hex8 FEA $\to$ 128-cube voxel grid.}

  \small
  The 3DWire dataset provides 14,293 building wireframe skeletons. Each is processed through
  a custom pipeline: wall/roof/floor solid extrusion, boolean union into one watertight body,
  and static FEA under ASCE~7-22 gravity + wind load combos using SfePy. Of 14,293 inputs,
  \textbf{3,115 (21.8\%)} are rejected for divergence, near-zero strain energy, or degenerate
  geometry, leaving \textbf{11,178 valid} $128^3$ voxel grids with 7 geometric channels and
  three scalar FEA labels: peak VM stress, max displacement, global compliance.
  A \textbf{family-aware split} groups wireframes by architectural cluster, preventing
  near-duplicate leakage: 8,943 train / 1,121 val / 1,114 test.
  Raw FEA outputs total $\sim$200\,GB; processed dataset is 59\,GB.
  Targets span 4.9 orders of magnitude (stress) to 7.7 orders (compliance), requiring
  log-transform normalization and 2nd/98th-percentile winsorization before training.
\end{minipage}%
\hspace{2mm}%
\begin{minipage}[t]{\HC}
  \sub{(2) Deep Ensemble Surrogate (5\,$\times$\,8.76M params)}
  \fig{\EX{Image 7.jpg}}{\linewidth}\\[-1pt]
  \fc{Fig.~3.\ Architecture: 4 conv stages (BN+GELU) $\to$ 3 SE-ResBlocks $\to$ dual pool $\to$ 3 outputs.}

  \small
  Each ensemble member is a 3D convolutional network.
  The architecture uses four strided conv stages (BatchNorm + GELU), followed by three
  Squeeze-and-Excitation ResBlocks for long-range spatial dependency, a dual global pooling
  head (avg + max), and three fully-connected layers down to 3 scalar outputs
  (stress, displacement, compliance). Five members are trained independently with different
  random seeds. The ensemble predicts \textbf{mean $\mu$} and \textbf{std.\ dev.\ $\sigma$}
  for each quantity, providing principled uncertainty bounds.
  Training: AdamW ($lr$=5e-4), Huber loss, cosine annealing, EMA ($\alpha$=0.999),
  random 90$^\circ$ rotations + reflections. Breakthrough at \textbf{model version~14}: first
  substantially accurate validation predictions (compliance Spearman $\rho$\,=\,0.948).
\end{minipage}

\vspace{3mm}\noindent\rule{\linewidth}{0.3pt}\vspace{2mm}

% SASTO Algorithm | 6-Connectivity
\begin{minipage}[t]{\HC}
  \sub{(3) SASTO Sensitivity-Guided Erosion}
  \fig{\EX{Image 8.jpg}}{\linewidth}\\[-1pt]
  \fc{Fig.~4.\ SASTO three-phase algorithm flowchart.}

  \small
  The surrogate replaces FEA inside a voxel-removal loop. Each iteration:
  \begin{itemize}[leftmargin=1.2em, itemsep=1pt, topsep=1pt]
    \item Compute distance transform; find candidates deeper than $t_{\min}(p)$.
    \item Score by ensemble sensitivity $s_i$; rank descending.
    \item Select batch $B$ of \textbf{6-simple-point} voxels (topology-safe by construction).
    \item Tentatively remove batch; query ensemble. If $\hat{\sigma}^+ \leq \sigma_\text{allow}$
      \textit{and} $\hat{C}/C_0 \leq 1.15$: \textbf{commit}. Else: \textbf{undo}
      $+$ halve $B$ (trust-region).
    \item \textbf{Phase~2} (endgame): fix $B$ at 5, then 1 for fine-grained pruning.
    \item \textbf{Phase~3}: swap moves, replacing corner-touching non-simple voxels with
      face-adjacent topology-safe alternatives.
  \end{itemize}
  Post-process: pocket fill $\to$ SDF smoothing $\to$ Marching Cubes $\to$ STL export.

  \eqpill{$\hat{\sigma}^+ = \mu_\sigma + k\,\sigma_\sigma$,\;
    $\hat{C}^+ = \mu_C + k\,\sigma_C$,\; $k=1.0$ at operating point.}
\end{minipage}%
\hspace{2mm}%
\begin{minipage}[t]{\HC}
  \sub{(4) Topology Preservation: 6-Connectivity}
  \begin{minipage}[c]{0.47\linewidth}\centering
    \fig{\EX{Image 9.jpg}}{13.3cm}\\[-1pt]
    \fc{{\color{posterred}26-conn: fragments}}
  \end{minipage}\hfill
  \begin{minipage}[c]{0.47\linewidth}\centering
    \fig{\EX{Image 10.jpg}}{13.3cm}\\[-1pt]
    \fc{{\color{teal}6-conn: 1 component}}
  \end{minipage}

  \vspace{1mm}\small
  \textit{Proposition (Kong \& Rosenfeld 1989): A binary voxel field with exactly one
  6-connected foreground component yields a single-component Marching Cubes mesh.}

  Two voxels are \textbf{6-adjacent} iff they share a face. They are
  \textbf{26-adjacent} iff they share a face, edge, or corner. Classical voxel removal
  using 26-adjacency produces isolated islands attached only at diagonal corners that
  become floating geometry after meshing, making them unfabricable. SASTO enforces the
  \textbf{6-simple-point} criterion: a voxel is removable only if deletion does not
  alter the 6-connected component count. This is checked in O(1) via a
  26-neighborhood Euler characteristic table lookup.

  \textbf{Experimental verification (60 designs):}
  \begin{itemize}[leftmargin=1.2em, itemsep=1pt, topsep=1pt]
    \item 100\% single-component voxel fields
    \item 90\% produce single-component meshes without post-processing
    \item 10\% resolved by a single connected-component fill pass
    \item Zero cases required geometry discard
  \end{itemize}

  \vspace{2mm}
  \sub{Training Curves \& Model Tuning}
  \begin{minipage}[c]{0.58\linewidth}
    \fig{\PF{fig_training_curves.png}}{\linewidth}\\[-1pt]
    \fc{Fig.~5.\ 5 ensemble members; convergence by epoch\,$\sim$130.}
  \end{minipage}\hfill
  \begin{minipage}[c]{0.15\linewidth}
    \fig{\PF{fig_activation.png}}{\linewidth}\\[-1pt]
    \fc{GELU}
  \end{minipage}\hfill
  \begin{minipage}[c]{0.22\linewidth}
    \small GELU outperforms ReLU on smooth stress fields. EMA reduces checkpoint noise. Removing SE blocks drops displacement $\rho$ from 0.970 to 0.921.
  \end{minipage}
\end{minipage}

\end{ccard}
\vspace{\SG}

%% C2: RESULTS & IN-SILICO VALIDATION ───────────────────────
\sechdr{Results \& In-Silico Validation}
\begin{ccard}

\begin{minipage}[t]{\TC}  %% COL A: Reference Case + Volume Distribution

  \sub{Reference Case (Sample 00472)}
  \begin{minipage}[c]{0.48\linewidth}\centering
    \fig{\RD{REF_original_colored_iso.png}}{8.4cm}\\[-1pt]\fc{Original}
  \end{minipage}\hfill
  \begin{minipage}[c]{0.48\linewidth}\centering
    \fig{\RD{ref_v11_pa_colored_iso.png}}{8.4cm}\\[-1pt]\fc{SASTO-PA ($-$45.0\%)}
  \end{minipage}
  \begin{minipage}[c]{0.48\linewidth}\centering
    \fig{\RD{REF_original_cutaway_iso.png}}{8.4cm}\\[-1pt]\fc{Cutaway: original}
  \end{minipage}\hfill
  \begin{minipage}[c]{0.48\linewidth}\centering
    \fig{\RD{ref_v11_pa_cutaway_iso.png}}{8.4cm}\\[-1pt]\fc{Cutaway: SASTO-PA}
  \end{minipage}

  \vspace{1mm}\small
  \begin{tabular}{@{}ll@{}}
    Volume removed: & $116{,}872 \to 64{,}292$ vox ($-$45.0\%) \\
    VM stress (cons.): & $3.08\!\times\!10^6$\,Pa \\
    Compliance ratio: & 1.004 (limit: 1.15) \\
    Mesh components: & 1 (watertight) \\
    Runtime: & 159.5\,s \\
    PA vs.\ U gain: & {\color{teal}\bfseries +10.7\,pp} \\
  \end{tabular}

  \vspace{3mm}
  \fig{\PF{fig_histogram.png}}{\linewidth}\\[-1pt]
  \fc{Fig.~6.\ Volume reduction distribution: 1,114 test geometries. Median 23.5\%.}

  \small
  Volume reductions span 5--65\%, reflecting the geometric diversity of 3DWire.
  SASTO-PA consistently outperforms SASTO-U by exploiting part-aware thickness floors
  for aggressive interior wall thinning. The negative-skew tail reflects complex
  houses with many load-bearing internal walls that resist further erosion.

\end{minipage}\hfill
%
\begin{minipage}[t]{\TC}  %% COL B: Multi-Geometry + Speedup

  \sub{Diverse Geometry Results}
  \fig{\HQ{fig_diverse_stl_gallery.png}}{\linewidth}\\[-1pt]
  \fc{Fig.~7.\ Four held-out test geometries: original (left), SASTO-PA solid (center), SASTO-PA cutaway (right). Reductions range from 18--41\%.}

  \small
  SASTO generalizes across markedly different house archetypes: simple rectangular plans,
  L-shaped footprints, multi-gabled roofs, and complex interior arrangements.
  All 1,114 test designs pass independent FEA re-analysis with zero violations.
  Volume reductions range from 18\% (geometry~08018, dense load-bearing interior walls)
  to 41\% (geometry~04203, large open-plan interior). The ``compressibility'' of a
  design is primarily determined by the fraction of interior partition surface area
  relative to total structure; open-plan designs with fewer interior walls allow
  more aggressive bulk exterior wall thinning without constraint violation.

  \vspace{2mm}
  \fig{\HQ{fig_type_comparison.png}}{\linewidth}\\[-1pt]
  \fc{Fig.~8.\ SASTO-U vs.\ SASTO-PA side-by-side. Part-aware removes more interior material while preserving the exterior structural envelope.}

  \vspace{2mm}
  \fig{\PF{fig_speedup.png}}{\linewidth}\\[-1pt]
  \fc{Fig.~9.\ Runtime comparison (log scale). SASTO 50\,s median vs.\ SIMP 19--77\,min projected.}

\end{minipage}\hfill
%
\begin{minipage}[t]{\TC}  %% COL C: FEA Validation + Cross-sections

  \sub{Independent FEA Validation ($n$\,=\,1,114)}
  \fig{\PF{fig_fea_compliance.png}}{\linewidth}\\[-1pt]
  \fc{Fig.~10.\ Compliance ratio $C_\text{opt}/C_\text{base}$ for all 1,114 test designs. Dashed = limit 1.15. Max observed: 1.004.}

  \small
  Every optimized design is re-analyzed by an \textit{independent} hex8 voxel FEA
  solver (not the trained surrogate) to provide a ground-truth safety check:
  \begin{itemize}[leftmargin=1.2em, itemsep=1pt, topsep=1pt]
    \item Mean $C_\text{opt}/C_\text{base} = 0.631 \pm 0.112$
    \item Maximum ratio: \textbf{1.004} (37\,pp below the 1.15 limit)
    \item {\color{teal}\bfseries Zero violations across all 1,114 designs}
    \item 99\% conformal upper bound: 0.950
  \end{itemize}

  \vspace{2mm}
  \fig{\HQ{fig_cross_section_comparison.png}}{\linewidth}\\[-1pt]
  \fc{Fig.~11.\ Cross-sections: original (top), SASTO-U (mid), SASTO-PA (bottom). Part-aware shows aggressive interior thinning.}

  \vspace{2mm}
  \begin{minipage}[c]{0.48\linewidth}
    \fig{\PF{fig_regression.png}}{\linewidth}\\[-1pt]\fc{Fig.~12.\ Regression}
  \end{minipage}\hfill
  \begin{minipage}[c]{0.48\linewidth}
    \fig{\PF{fig_bland_altman.png}}{\linewidth}\\[-1pt]\fc{Fig.~13.\ Bland-Altman}
  \end{minipage}

  \small\vspace{1mm}
  The regression and Bland-Altman plots confirm that for designs where the surrogate
  predicts low compliance (efficient designs), the FEA ground truth also records low
  compliance, validating the surrogate's ranking accuracy as the mechanism underlying
  SASTO's safety record.

\end{minipage}

\vspace{3mm}
\eqpill{Conservative bound: $\hat{\sigma}^+ = \mu + k\sigma$ ($k$=1.0).
  $P(\text{violation}) \leq \tfrac{1}{n+1} = 0.09\%$
  (Vovk et al.\ conformal guarantee, distribution-free, $n$=1,114).
  Safety margin to compliance limit: 0.146 units.}
{\centering\tiny\itshape
  Model v14; 5-member deep ensemble; independent hex8 FEA validation on all 1,114 test designs.\par}

\end{ccard}

\end{minipage}%  ─── END CENTER PANEL ───
\hspace{3mm}%
%% ════════════════════════════════════════════════════════════
%%  RIGHT PANEL  (29.72cm)
%% ════════════════════════════════════════════════════════════
\begin{minipage}[t]{\RC}

%% R1: STATISTICAL ANALYSIS ────────────────────────────────
\sechdr{Statistical Analysis}
\begin{ccard}
  \small\RaggedRight
  I conducted a comprehensive statistical analysis to characterize SASTO's performance
  across the full 1,114-geometry test set, including surrogate quality, optimization
  outcomes, uncertainty calibration, and runtime distribution.

  \vspace{3mm}
  \sub{k-Factor Pareto Frontier}
  \fig{\PF{fig_k_factor.png}}{\linewidth}\\[-1pt]
  \fc{Fig.~14.\ $k$-factor: acceptance rate (teal, left) and mean volume reduction (gold, right). Operating point $k$=1.0.}

  \small
  The scalar $k$ controls the safety margin in the conservative bound $\hat{y}^+ = \mu + k\sigma$.
  At $k$=0 the surrogate mean is used directly; at $k$=3 the bound is extremely conservative.
  The dual-axis Pareto plot reveals that larger $k$ reduces per-design material savings but
  provides stronger safety guarantees. \textbf{$k$=1.0} was selected as the operating point:
  it maintains $>$99\% acceptance rate while providing a $1\sigma$ buffer against
  surrogate under-prediction, corresponding to the observed conformal $k$=1.90 for 84.1\%
  compliance coverage.

  \vspace{3mm}
  \sub{Ensemble Uncertainty Bands}
  \fig{\PF{fig_uncertainty.png}}{\linewidth}\\[-1pt]
  \fc{Fig.~15.\ Ensemble mean $\pm\sigma$ during reference-case optimization. Bands narrow as design converges.}

  \small
  Ensemble uncertainty $\sigma$ naturally decreases as optimization converges toward a
  structurally efficient configuration, because such designs are well-represented in
  the training distribution. This self-calibrating behavior provides \textit{additional}
  implicit safety margin precisely when the design is most aggressively pruned.

  \vspace{3mm}
  \sub{Surrogate Quality (test set, $n$\,=\,1,114)}
  {\small
  \begin{tabular}{@{}lrrr@{}}
    \rowcolor{secbar}
    \color{white}\bfseries Target &
    \color{white}\bfseries Spear.~$\rho$ &
    \color{white}\bfseries $R^2_{\log}$ &
    \color{white}\bfseries MAPE \\\midrule
    VM stress    & 0.737 & 0.419 & 37.4\% \\
    \rowcolor{altrow}
    Displacement & 0.970 & 0.842 & 10.9\% \\
    Compliance   & 0.948 & 0.814 & 18.5\% \\
  \end{tabular}}

  \vspace{1mm}\small
  The surrogate does not need pointwise accuracy; it needs to correctly
  \textit{rank} designs by compliance so that sensitivity scores steer erosion away
  from essential voxels. Compliance $\rho$=0.948 far exceeds the $\geq$0.90 endpoint.
  VM stress MAPE is higher (37.4\%) due to localized stress concentrations, but the
  $k$-bound compensates: $k$=1.0 adds ${\sim}1\sigma\approx 0.37$ compliance units of
  margin, which is 2.5$\times$ the observed reserve (0.146 units).

  \vspace{3mm}
  \sub{Conformal Prediction Certificates}
  \begin{itemize}[leftmargin=1.2em, itemsep=1pt, topsep=0pt]
    \item $k$ for 84.1\% compliance coverage: $k=1.90$ (heavier tails than Gaussian)
    \item $k$ for 84.1\% VM stress coverage: $k=4.31$ (localized peaks)
    \item 99\% conformal upper bound on $C_\text{opt}/C_\text{base}$: \textbf{0.950}
    \item {\color{teal}\bfseries $P(\text{viol.})\leq 1/(n{+}1)=0.09\%$, distribution-free}
  \end{itemize}

  \vspace{3mm}
  \sub{Runtime Summary}
  {\small
  \begin{tabular}{@{}lll@{}}
    \rowcolor{secbar}
    \color{white}\bfseries Method &
    \color{white}\bfseries Per Design &
    \color{white}\bfseries vs.\ SASTO \\\midrule
    SASTO (ours)  & 50\,s median & baseline \\
    \rowcolor{altrow}
    SIMP 64-cube  & 94\,s median & 0.5$\times$ \\
    SIMP 128-cube & 19--77\,min  & 23--92$\times$ slower \\
  \end{tabular}}

  \small\vspace{1mm}
  Measured on RTX A3000 Laptop GPU. SIMP 128-cube times are projected from measured
  64-cube data with $O(N^2)$ scaling (each additional resolution doubling $\approx 4\times$
  solve time). SASTO's 50\,s runtime is dominated by ensemble inference ($\sim$80\%) and
  distance-transform ($\sim$15\%), both highly GPU-parallelizable.
\end{ccard}
\vspace{\SG}

%% R2: CONCLUSIONS ─────────────────────────────────────────
\sechdr{Conclusions}
\begin{ccard}
  \small\RaggedRight
  This research demonstrates that a 5-member deep ensemble surrogate can replace iterative
  FEA within a topology optimization loop for full building-scale structures. SASTO achieves
  competitive material savings at a fraction of the compute cost, while guaranteeing
  printable single-component topology and distribution-free safety coverage. Specifically,
  SASTO solves the three identified failure modes: it is fast (50\,s), enforces 6-connected
  topology for printability, and provides conformal safety certificates.\\[4pt]
  \textbf{Revisiting Engineering Criteria:}
  \begin{itemize}[leftmargin=1.2em, itemsep=1pt, topsep=2pt]
    \item \textbf{Speed:} 50\,s median vs.\ SIMP 19--77\,min at equivalent resolution.
      \textbf{Endpoint: met.}
    \item \textbf{Printability:} 100\% single-component voxel fields; 90\% direct
      single-component meshes. \textbf{Endpoint: met.}
    \item \textbf{Safety:} 0/1,114 FEA violations; 99\%-conformal bound 0.950 vs.\
      limit 1.15. \textbf{Endpoint: met and exceeded.}
  \end{itemize}

  \vspace{3mm}
  \textbf{Scientific Contributions.}
  SASTO is the first system combining (a)~building-scale 3D-CNN surrogate inference,
  (b)~6-connected topology enforcement for printability, and (c)~conformal safety
  coverage for structural topology optimization. The 6-connectivity proposition and
  family-aware data splitting are novel contributions applicable beyond this domain.\\[3pt]
  \textbf{Practical Impact.}
  At scale, 23.5\% concrete reduction per building translates to millions of tonnes of
  avoided CO$_2$ annually. SASTO runs on a single consumer GPU in under one minute,
  enabling engineers to explore dozens of design variants in a morning, a workflow
  that was previously impossible with FEA-based optimization.
  {\color{teal}\bfseries 8\% of global CO$_2$ is cement. Even 10\% adoption of
  AI-optimized construction would offset tens of millions of tonnes annually.}
\end{ccard}
\vspace{\SG}

%% R3: FUTURE WORK ─────────────────────────────────────────
\sechdr{Future Work}
\begin{ccard}
  \small\RaggedRight
  This work opens several promising research directions, distinguished by near-term
  feasibility vs.\ requiring additional data or infrastructure:

  \vspace{2mm}
  \begin{itemize}[leftmargin=1.2em, itemsep=2pt, topsep=0pt]
    \item \textbf{FEA-in-the-loop active learning:} When ensemble $\sigma$ exceeds a
      threshold mid-run, trigger a ground-truth FEA solve, fine-tune the member on the
      new data point, and continue, providing a self-correcting safety net.
    \item \textbf{Nonlinear material spot checks:} Apply concrete damaged plasticity
      (CDP) FEA to assess cracking and buckling in 78\,mm interior partitions under
      both service and extreme loads.
    \item \textbf{Physical print validation:} 3D-print one optimized house at 1:10 scale
      in micro-concrete; apply incremental compression load to failure; compare DIC
      strain fields to FEA predictions.
    \item \textbf{Seismic + wind} load extensions for code compliance in ASCE~7-22
      Chapter~12 seismically active zones.
    \item \textbf{Multi-material:} Allow geopolymer or fiber-reinforced concrete in thin
      sections for a higher strength-to-weight ratio without increasing volume.
    \item \textbf{Other domains:} Ship hulls, aircraft ribs, biomedical scaffolds; all share the same boundary-value problem structure.
  \end{itemize}
\end{ccard}
\vspace{\SG}

%% R4: KEY REFERENCES ──────────────────────────────────────
\sechdr{Key References}
\begin{ccard}
  \footnotesize
  \begin{enumerate}[leftmargin=1.4em, itemsep=2pt, topsep=0pt, label={[\arabic*]}]
    \item M.~P.~Bendsoe \& O.~Sigmund. \textit{Topology Optimization.} Springer, 2003.
    \item B.~Lakshminarayanan et al. Deep ensembles for uncertainty.
      \textit{NeurIPS}, 2017.
    \item T.~Y.~Kong \& A.~Rosenfeld. Digital topology survey.
      \textit{CVGIP} 48(3):357--393, 1989.
    \item ASCE. \textit{ASCE/SEI~7-22: Minimum Design Loads.} 2022.
    \item R.~A.~Buswell et al. 3D printing using concrete extrusion.
      \textit{Cem.\ Concr.\ Res.} 112:37--49, 2018.
    \item O.~Sigmund \& K.~Maute. Topology optimization approaches.
      \textit{Struct.\ Multidisc.\ Optim.} 48:1031, 2013.
    \item Y.~Lin et al. 3DWire wireframe dataset. KAUST, 2024.
    \item IEA. \textit{Global Buildings Report 2021.} Int'l Energy Agency.
    \item W.~E.~Lorensen \& H.~E.~Cline. Marching Cubes.
      \textit{ACM SIGGRAPH} 21(4):163--169, 1987.
    \item V.~Vovk et al. \textit{Algorithmic Learning in a Random World.}
      Springer, 2005.
  \end{enumerate}
  \vspace{2mm}
  {\centering\color{posterred}\itshape Full bibliography in accompanying paper.\par}
\end{ccard}

\end{minipage}%  ─── END RIGHT PANEL ───

\end{document}
